% IEEE two-column report. Figures and tables are read from ../results (see the \figdir and \tabdir macros).
% Build: latexmk -pdf main.tex   (needs a LaTeX install with IEEEtran and booktabs)
\documentclass[conference]{IEEEtran}
\usepackage{graphicx}
\usepackage{booktabs}
\usepackage{amsmath}
\usepackage{amssymb}
\usepackage{array}
\usepackage{xcolor}
\usepackage[hidelinks]{hyperref}
\newcommand{\TODO}[1]{\textcolor{red}{[TODO: #1]}}
\newcommand{\figdir}{../results/figures}
\newcommand{\tabdir}{../results/tables}

\begin{document}

\title{Label-Efficient Land-Use Classification on EuroSAT:\\
RGB versus 13-Band Input, Supervised versus Pseudo-Labelled Training}

\author{\IEEEauthorblockN{Ishant Bhoyar}
\IEEEauthorblockA{Enrolment no. 202410101200059\\Deep Learning, Sitare University}
\and
\IEEEauthorblockN{Gaurav Kumar}
\IEEEauthorblockA{Enrolment no. 202410101200047\\Deep Learning, Sitare University}
\and
\IEEEauthorblockN{Jagriti Pandey}
\IEEEauthorblockA{Enrolment no. 202410101200060\\Deep Learning, Sitare University}
\and
\IEEEauthorblockN{Ayush Kumar}
\IEEEauthorblockA{Enrolment no. 202410101200028\\Deep Learning, Sitare University}}

\maketitle

\begin{abstract}
We ask how much test accuracy a land-use classifier loses when only 1\% or 10\% of the EuroSAT training labels are used, whether the extra Sentinel-2 bands help, and whether pseudo-labelling recovers some of the loss. We train a small CNN from scratch, an ImageNet-pretrained ResNet-50, and an ImageNet-pretrained ViT-Small on RGB and on all 13 bands, at 1\%, 10\% and 100\% of the labels (three seeds at 1\% and 10\%, one at 100\%). With 1\% of the labels (191 images), the pretrained models reach 88.5--91.2\% test macro-F1 and the scratch CNN 74.6--78.1\%; with all labels every model reaches 97.5--99.1\%. The 13-band input helps the scratch CNN most (+3.5 to +4.2 points at 1\% and 10\%), ResNet-50 by about +0.9 points, and ViT-Small by 0.5 points or less, which at 1\% is within the spread between seeds. A control from an independent pipeline of a teammate shows that RGB cut from the 13-band files is itself about 1 point better than the JPEG RGB we used as baseline, which explains the ResNet-50 gain at 1\% of the labels. Pseudo-labelling a ResNet-50 improves it by 2.3 to 2.4 points at 1\% of the labels and by 0.2 points at 10\%. Pseudo-labelling results rest on two seeds, and the 100\% results on one.
\end{abstract}

\begin{IEEEkeywords}
EuroSAT, Sentinel-2, label efficiency, multispectral, pseudo-labelling, transfer learning
\end{IEEEkeywords}

\section{Introduction}
Land-use and land-cover maps built from satellite images need labelled examples, and labels are expensive: an expert has to inspect each patch. Sentinel-2 satellites, however, record far more than the red, green and blue (RGB) bands of an ordinary photograph. They measure 13 spectral bands, including near-infrared and short-wave infrared, which separate vegetation, water and built-up areas better than visible light does~\cite{drusch2012sentinel}. Two ideas could therefore reduce the labelling effort: using the extra spectral bands, and using the many unlabelled images that are available anyway.

We study these ideas on the EuroSAT benchmark~\cite{helber2019eurosat} and ask three questions.
\begin{enumerate}
\item \textbf{Label efficiency.} How quickly does test accuracy fall when only 10\% or 1\% of the training labels are available, for a small network trained from scratch, an ImageNet-pretrained convolutional network, and an ImageNet-pretrained vision transformer?
\item \textbf{Spectral bands.} Does giving the same models all 13 bands, instead of RGB, help, and is the gain larger when labels are scarce?
\item \textbf{Unlabelled data.} Does pseudo-labelling, which adds confident predictions on unlabelled training images as extra training labels, improve a supervised model trained on the same few labels?
\end{enumerate}
Every comparison uses the same fixed split, the same training budget rule and three random label subsets. Our contributions are: (i) a reproducible protocol with nested label subsets, leakage checks and a train-only normalisation; (ii) a way to start a pretrained RGB network from its ImageNet weights when 13 bands are given; (iii) a pseudo-labelling procedure whose confidence threshold is chosen on the validation set only; and (iv) paired comparisons matched on the seed, so that two settings are always compared on the same labelled subset.

The rest of the report is organised as follows. Section~\ref{sec:related} places the work in the literature. Section~\ref{sec:data} describes the data and the splits, Section~\ref{sec:method} the models, the training procedure and the pseudo-labelling method, and Section~\ref{sec:setup} the experimental setup, including what was cut to fit the time limit and where the numbers come from. Section~\ref{sec:results} reports the results, Section~\ref{sec:discussion} answers the three questions and gives practical recommendations, and Sections~\ref{sec:limits} and~\ref{sec:future} list the limitations and future work. The appendices list every run, the design decisions and the steps to reproduce the results.

\section{Related Work}\label{sec:related}
\textbf{Data and benchmarks.} Sentinel-2 is a constellation of optical satellites that image the land surface in 13 spectral bands at 10, 20 and 60~m resolution~\cite{drusch2012sentinel}. EuroSAT~\cite{helber2019eurosat} provides 27{,}000 labelled Sentinel-2 patches in 10 classes, in RGB and in all 13 bands. The original paper trained convolutional networks on it and reported an overall accuracy of about 98.6\% for ResNet-50 on RGB with a large labelled training set; we use this as a sanity check for our own baseline. BigEarthNet~\cite{sumbul2019bigearthnet} is a much larger Sentinel-2 archive with multiple labels per patch. Reviews of deep learning in remote sensing~\cite{zhu2017review} list label scarcity as a main obstacle.

\textbf{Transfer learning.} Convolutional networks pretrained on ImageNet~\cite{he2016resnet} and vision transformers~\cite{dosovitskiy2021vit} are the usual starting points when labels are few. How well features transfer depends on how close the source and target data are~\cite{yosinski2014transfer}, and for remote sensing, representations learned on in-domain imagery have been reported to transfer better than ImageNet ones~\cite{neumann2019indomain}. Recent work pretrains on unlabelled multispectral Sentinel data with self-supervision~\cite{wang2023ssl4eo,cong2022satmae}. We do not pretrain on multispectral data; we adapt RGB-pretrained weights to 13 bands (Section~\ref{sec:models}), because that is the only option when a multispectral checkpoint is not available. Both ImageNet families expect three input channels, so using more than three bands requires changing the first layer.

\textbf{Pretrained checkpoints.} Our pretrained models come from the \texttt{timm} library~\cite{wightman2019timm}. The ResNet-50 weights follow the improved training recipe of Wightman et al.~\cite{wightman2021rsb}, and the ViT-Small weights follow the augmentation and regularisation study of Steiner et al.~\cite{steiner2022augreg}. Both were trained at $224\times224$, which is why we upsample the $64\times64$ patches.

\textbf{Semi-supervised learning.} Pseudo-labelling~\cite{lee2013pseudo} trains a model on its own confident predictions for unlabelled images. Self-training with a noisy student~\cite{xie2020noisystudent} trains an equal or larger student, with added noise, on the labelled data plus the teacher's labels. MixMatch~\cite{berthelot2019mixmatch} and FixMatch~\cite{sohn2020fixmatch} add consistency losses between differently augmented views of each unlabelled image. Oliver et al.~\cite{oliver2018realistic} stress that semi-supervised methods must be compared with a well-tuned supervised baseline on the same budget, and Arazo et al.~\cite{arazo2020confirmation} show that wrong pseudo-labels can reinforce themselves (confirmation bias). We use the simple pseudo-labelling form, a strict choice of the threshold on validation data only, and a comparison of every student with its own teacher; consistency methods are left as future work.

\section{Data and Splits}\label{sec:data}
EuroSAT has 27{,}000 patches of $64\times64$ pixels in 10 classes, with 2{,}000 to 3{,}000 images per class (Table~\ref{tab:classes}). We use two versions of it: an RGB version (8-bit JPEG renderings of bands B04, B03, B02) and a 13-band version (16-bit GeoTIFF). Both contain exactly the same patch identifiers, so one split serves both inputs and an RGB model and a 13-band model are always trained and tested on the same scenes.

\begin{table}[t]
\centering
\caption{Images per class and split (stratified 70/15/15 split, seed 42).}
\label{tab:classes}
\small
\begin{tabular}{lrrrr}
\toprule
Class & Train & Val. & Test & Total \\
\midrule
AnnualCrop & 2{,}100 & 450 & 450 & 3{,}000 \\
Forest & 2{,}100 & 450 & 450 & 3{,}000 \\
HerbaceousVegetation & 2{,}099 & 451 & 450 & 3{,}000 \\
Highway & 1{,}750 & 375 & 375 & 2{,}500 \\
Industrial & 1{,}750 & 375 & 375 & 2{,}500 \\
Pasture & 1{,}400 & 300 & 300 & 2{,}000 \\
PermanentCrop & 1{,}750 & 375 & 375 & 2{,}500 \\
Residential & 2{,}100 & 450 & 450 & 3{,}000 \\
River & 1{,}750 & 375 & 375 & 2{,}500 \\
SeaLake & 2{,}100 & 450 & 450 & 3{,}000 \\
\midrule
Total & 18{,}899 & 4{,}051 & 4{,}050 & 27{,}000 \\
\bottomrule
\end{tabular}
\end{table}

\textbf{Splits.} We split the data once, stratified by class, with a fixed seed: 70\% train (18{,}899 images), 15\% validation (4{,}051) and 15\% test (4{,}050). The validation and test sets are identical for every model, input type, label fraction and seed. The validation set is used for early stopping, model selection, learning-rate choice and the pseudo-label threshold; the test set is used only to report the final scores.

\textbf{Label fractions.} From the training split we draw labelled subsets of 1\% (191 images), 10\% (1{,}890) and 100\% of the images, each with three seeds. Within a seed the subsets are nested ($1\% \subset 10\% \subset 100\%$), so adding labels never removes examples already in use, and each subset keeps the class proportions. The remaining training images form the \emph{unlabelled pool} used for pseudo-labelling (18{,}708 images at 1\%, 17{,}009 at 10\%); their labels are never read during training. Validation and test images are never used as unlabelled data. Scripts check for overlaps and store hashes of every split file, and training stops if a file has been changed.

\textbf{Preprocessing.} Each channel is standardised with the mean and standard deviation of the training split only. Augmentation uses the same operations for 3 and 13 channels: random flips and $90^\circ$ rotations, and a random resized crop. Images stay at $64\times64$; models that need $224\times224$ input upsample on the GPU with bilinear interpolation.

\begin{table}[t]
\centering
\caption{The 13 Sentinel-2 bands in the order of the EuroSAT 13-band files. Channel 12 holds B8A.}
\label{tab:bands}
\small
\begin{tabular}{rlrrl}
\toprule
Ch. & Band & Centre (nm) & Res. (m) & Note \\
\midrule
0 & B01 & 443 & 60 & coastal aerosol \\
1 & B02 & 490 & 10 & blue \\
2 & B03 & 560 & 10 & green \\
3 & B04 & 665 & 10 & red \\
4 & B05 & 705 & 20 & red edge \\
5 & B06 & 740 & 20 & red edge \\
6 & B07 & 783 & 20 & red edge \\
7 & B08 & 842 & 10 & near-infrared \\
8 & B09 & 945 & 60 & water vapour \\
9 & B10 & 1375 & 60 & cirrus \\
10 & B11 & 1610 & 20 & SWIR \\
11 & B12 & 2190 & 20 & SWIR \\
12 & B8A & 865 & 20 & narrow near-infrared \\
\bottomrule
\end{tabular}
\end{table}

\textbf{Band order.} The 13-band files carry no band names, and many code bases assume the standard Sentinel-2 order (B8A between B08 and B09). We found that band B8A is stored \emph{last} (order B01--B12, then B8A; Table~\ref{tab:bands}) by checking that channels 3, 2, 1 correlate above 0.99 with the RGB channels, that channel 9 has the very low mean of cirrus band B10, and that channels 10 and 11 match the short-wave infrared bands. With the standard order, five bands would get the wrong name. The band order matters for the first-layer initialisation described below, which copies the RGB filters to channels 3, 2 and 1.

\section{Method}\label{sec:method}
\subsection{Models}
\label{sec:models}
We compare three model families (Table~\ref{tab:models}).
\begin{itemize}
\item \textbf{Scratch CNN} (2.8M parameters): a small ResNet-style network with four stages of two basic blocks (widths 32--256), trained from random initialisation on the $64\times64$ input.
\item \textbf{ResNet-50}~\cite{he2016resnet} (23.5M): the ImageNet-1k weights \texttt{resnet50.a1\_in1k} from the \texttt{timm} library, fed images upsampled to $224\times224$.
\item \textbf{ViT-Small}~\cite{dosovitskiy2021vit} (21.7M): the \texttt{timm} model \texttt{vit\_small\_patch16\_224.}\allowbreak\texttt{augreg\_in21k\_ft\_in1k}, pretrained on ImageNet-21k and fine-tuned on ImageNet-1k, also at $224\times224$.
\end{itemize}

\begin{table}[t]
\centering
\caption{The three models. Parameters are given for RGB and for 13-band input.}
\label{tab:models}
\small
\begin{tabular}{llcc}
\toprule
Model & Pretraining & Input & Params (RGB / 13) \\
\midrule
Scratch CNN & none & 64 px & 2.80M / 2.80M \\
ResNet-50 & ImageNet-1k & 224 px & 23.5M / 23.6M \\
ViT-Small & ImageNet-21k, then 1k & 224 px & 21.7M / 22.7M \\
\bottomrule
\end{tabular}
\end{table}

\textbf{Using 13 bands with pretrained weights.} For 13-band input we replace the first layer (the first convolution of ResNet-50, the patch-embedding projection of ViT-Small). The red, green and blue filters of the pretrained layer are copied to the matching bands B04, B03 and B02. Each of the other ten bands receives the mean of the three RGB filters. All 13 filters are then multiplied by $3/13$, so that the summed first-layer response keeps about the size the network saw with three channels (every band is standardised to unit variance). The scratch CNN simply has 13 input channels. This initialisation is a design choice with limits: ImageNet pretraining knows nothing about near-infrared or short-wave-infrared bands, so the network must learn how to use them during fine-tuning, which is hard with only 191 labelled images. The scaling by $3/13$ also weakens the RGB filters, so the first forward pass of the 13-band model differs from that of the RGB model.

\subsection{Training}
All runs use AdamW~\cite{loshchilov2019adamw} without weight decay on biases and normalisation layers, a linear warm-up~\cite{goyal2017warmup} followed by cosine decay~\cite{loshchilov2017sgdr}, cross-entropy loss, mixed-precision arithmetic~\cite{micikevicius2018mixed}, and batches of 64. The implementation uses PyTorch~\cite{paszke2019pytorch}. After every epoch the model is scored on the validation set; the weights with the best validation macro-F1 are kept, and training stops after 10 epochs without improvement (15 for the scratch CNN, which trains for up to 60 epochs). The test set is predicted once, with the selected weights. Table~\ref{tab:hparams} lists the hyper-parameters.

\begin{table}[t]
\centering
\caption{Hyper-parameters. Every run trains for $\max(\text{epochs}\times\text{one pass},\ \text{min steps})$ steps.}
\label{tab:hparams}
\small
\begin{tabular}{lccccc}
\toprule
Model & LR & Weight decay & Warm-up & Epochs & Min steps \\
\midrule
Scratch CNN & $2\times10^{-3}$ & 0.05 & 5\% & 60 & 3{,}000 \\
ResNet-50 & $5\times10^{-4}$ & 0.05 & 5\% & 30 & 1{,}500 \\
ViT-Small & $10^{-4}$ & 0.05 & 10\% & 30 & 1{,}500 \\
\bottomrule
\end{tabular}
\end{table}

\textbf{Equal training budget.} Every run trains for $\max(\text{epochs}\times\text{one pass over its labelled set},\ \text{min\_steps})$ optimiser steps, with 30 epochs for the pretrained models and a floor of 1{,}500 steps (3{,}000 for the scratch CNN). With 100\% of the labels one pass is 295 steps. With 1\% or 10\% of the labels one pass is only 2 or 29 steps, so each epoch is stretched to 50 steps, drawn with replacement. The number of epochs, and so the number of validation checks, is therefore the same for every label fraction. A consequence is that small label sets are passed over many more times than large ones, which can favour overfitting but also gives every fraction enough steps to converge.

\textbf{Learning rate.} We swept three learning rates on one cell (RGB, 10\% labels, seed 0), choosing by validation macro-F1 only; the test set was never evaluated (Table~\ref{tab:lr_sweep}). For the scratch CNN, $2\times10^{-3}$ was best among $\{10^{-3}, 2\times10^{-3}, 4\times10^{-3}\}$. For ResNet-50 the validation macro-F1 was 0.964, 0.967 and 0.969 for $2\times10^{-4}$, $5\times10^{-4}$ and $10^{-3}$; the differences are below 0.5 points, so we kept the default $5\times10^{-4}$. ViT-Small used its default $10^{-4}$ without a sweep. The chosen rate is used unchanged for every input type, label fraction and seed.

\input{\tabdir/lr_sweep.tex}

\subsection{Pseudo-labelling}
A supervised run with 1\% or 10\% of the labels acts as the \emph{teacher}. After training it predicts the unlabelled pool. The procedure for one teacher is:
\begin{enumerate}
\item Predict the class probabilities of every pool image with the teacher.
\item On the validation set, find the smallest threshold $\tau\in\{0.50,0.51,\dots,0.99\}$ at which the teacher's predictions with confidence $\geq\tau$ are at least 95\% correct. If no $\tau$ reaches 95\%, use the one with the highest precision.
\item Give every pool image with confidence $\geq\tau$ the teacher's class as its pseudo-label.
\item Train a \emph{student}, a fresh model with the same architecture, starting weights and hyper-parameters, on the labelled images plus the pseudo-labelled images under the same budget rule, with early stopping on the true validation labels.
\item Compare the student with its own teacher (same model, input, fraction and seed).
\end{enumerate}
The true labels of pool images are read only after training, to measure how accurate the pseudo-labels were; nothing in training or selection uses them. Choosing $\tau$ on the validation set means that no test data is used and that as many pseudo-labels as possible are kept at a known quality. Starting the student from fresh weights, and not from the teacher, keeps the comparison between the two like for like.

\subsection{Evaluation}
The main metric is test macro-F1, the mean over the $K=10$ classes of $F1_k = 2P_kR_k/(P_k+R_k)$ with the precision $P_k$ and recall $R_k$ of class $k$; it gives every class the same weight, which suits mildly imbalanced classes. Test accuracy is also reported. Settings are compared with paired differences $\Delta = F1_b - F1_a$ matched on seed, so that both runs saw the same labelled subset. Where saved predictions exist, the code (\texttt{scripts/evaluate.py}) also computes per-class F1, expected calibration error and 95\% paired bootstrap intervals~\cite{efron1979bootstrap}: the test images are resampled 1{,}000 times with replacement, the same images are used for both runs, and the interval is the 2.5th to 97.5th percentile of the resampled $\Delta$. Section~\ref{sec:logs} says which results have these.

\section{Experimental Setup}\label{sec:setup}
All runs were trained on free-tier cloud GPUs (NVIDIA Tesla T4, 15 GB). As a sanity check, ResNet-50 on RGB with 100\% of the labels reaches 98.67\% test accuracy and 98.63\% macro-F1 (43 minutes), in line with the about 98.6\% reported for this benchmark~\cite{helber2019eurosat}. A 1\% or 10\% run of a pretrained model takes about 10 minutes.

To fit the time limit the full grid was reduced in four ways. (i)~The 1\% and 10\% cells use three seeds (0, 1, 2), but the 100\% cells use one seed (0), so they have no standard deviation; this gives $3\times2\times2\times3 + 3\times2 = 42$ supervised runs. (ii)~Learning rates come from a short sweep on the validation set for ResNet-50 and the scratch CNN at 10\% labels (see the Method section); ViT-Small uses its default rate. (iii)~Pseudo-labelling was planned for ResNet-50 and the scratch CNN at 1\% and 10\% labels, for both inputs and three seeds, and not for ViT-Small; ResNet-50 students train for 10 epochs, about the same compute as their teachers. The training sessions ended before all of these finished, so only nine students exist: eight ResNet-50 students (both inputs, both fractions, seeds 0 and 1) and one scratch CNN student (RGB, 1\%, seed 0). (iv)~The optional experiments (class-balanced students, a second self-training round, a band ablation, FixMatch) were not run. The 51 finished runs took about 9.2 hours of T4 time, and the learning-rate sweeps about 40 minutes more. Software on the training machines: Python 3.13, PyTorch 2.14 (CUDA 13.0), timm 1.0.31. Test accuracy and macro-F1 of the lanes run on Colab were read from the printed logs of each run (four decimals), because only one lane kept the prediction files; Section~\ref{sec:logs} states what that changes.

\subsection{Compute cost}
Table~\ref{tab:training_cost} gives the mean time per supervised run. The equal-budget rule makes small-label runs cheap: a ResNet-50 run takes 8.5 minutes at 1\% of the labels and 42.4 minutes at 100\%, about five times longer. The scratch CNN is the cheapest model at every fraction, and ViT-Small is cheaper than ResNet-50 on the same GPU (33.1 against 42.4 minutes at 100\%). A pseudo-label student of ResNet-50 (10 epochs over the labelled and pseudo-labelled images) took 12.4 to 15.0 minutes, about 1.5 times the time of its teacher.

\input{\tabdir/training_cost.tex}

\subsection{Source of the numbers}\label{sec:logs}
Only the run folders of one lane (13 runs) were kept. For all runs, the test accuracy and macro-F1 are the values printed when each run finished; for the 13 kept runs they were checked against the saved metrics, and they agree to the four printed decimals. This has four consequences. (a)~Differences between two settings are per-seed differences of two macro-F1 values, without bootstrap intervals; only the ViT-Small RGB-versus-13-band comparison has intervals. (b)~Calibration error is not reported. (c)~Per-class scores and confusion matrices exist only for ViT-Small at 1\% and 10\% labels and ResNet-50 on 13 bands at 100\%. (d)~We call a difference consistent when it has the same sign in every seed; with two or three seeds this is a weak statement, and we avoid p-values.

\section{Results}\label{sec:results}
\subsection{Supervised label efficiency}
\begin{figure*}[t]
\centering
\includegraphics[width=0.85\textwidth]{\figdir/label_efficiency_f1.png}
\caption{Test macro-F1 against the share of labelled training images (mean $\pm$ std over three seeds at 1\% and 10\%; one seed at 100\%). The pretrained models lose 8--10 points from 100\% to 1\% of the labels; the scratch CNN loses 21--23.}
\label{fig:label_efficiency}
\end{figure*}
\input{\tabdir/supervised_f1.tex}
With all labels, every model is within 97.5--99.1\% macro-F1, and the differences between the pretrained models (at most 0.4 points on one seed) are not meaningful. Accuracy falls much more for the scratch CNN. From 100\% to 1\% of the labels, ResNet-50 loses 10.1 points on RGB and 9.7 on 13 bands, and ViT-Small 8.1 and 7.7; the scratch CNN loses 22.9 and 20.7. Test accuracy (Table~\ref{tab:supervised_acc}) is between 0.0 and 1.0 point higher than macro-F1 in every cell and picks the same best setting in every column, so the conclusions do not depend on the metric.

\input{\tabdir/supervised_acc.tex}

\subsection{Comparison of the models}
Table~\ref{tab:model_comparison} compares the models directly, matched on seed. At 1\% of the labels, ResNet-50 beats the scratch CNN by 13.9 (RGB) and 11.3 (13 bands) points, and ViT-Small beats it by 16.1 and 13.1 points, in all three seeds. ViT-Small is ahead of ResNet-50 at 1\% in all three seeds (+2.2 RGB, +1.8 13 bands) and at 10\% by 1.0 point on RGB (3/3 seeds); on 13 bands at 10\% the gap is 0.3 points and only 2 of 3 seeds agree. The gap between pretrained and scratch models shrinks with more labels and with the 13-band input: at 10\% of the labels ViT-Small is 5.7 points ahead of the scratch CNN on RGB but only 1.6 on 13 bands. At 100\% all gaps are 1.3 points or less on a single seed, and ResNet-50 on 13 bands (99.05\%) and ViT-Small (98.91\%) differ by 0.14 points, which we do not read as a difference.

\input{\tabdir/model_comparison.tex}

\subsection{Value of the extra spectral bands}
\begin{figure}[t]
\centering
\includegraphics[width=\columnwidth]{\figdir/rgb_vs_ms_gain.png}
\caption{Gain in test macro-F1 of 13-band over RGB input; bars are the mean over seeds, dots single seeds. The gain is largest for the scratch CNN and small for the pretrained models.}
\label{fig:gain}
\end{figure}
Thirteen bands improve every model at every fraction in the mean, but by very different amounts (Fig.~\ref{fig:gain}). The scratch CNN gains 3.5 points at 1\% and 4.2 at 10\% labels, in all three seeds, and 1.2 points at 100\% on one seed. ResNet-50 gains 0.9 points at 1\% and at 10\% (3/3 seeds each, standard deviation 0.6 and 0.07) and 0.4 at 100\%. ViT-Small gains 0.5 points at 1\% (2 of 3 seeds; the seed-to-seed standard deviation is 0.9) and 0.1 points at 10\%; the bootstrap interval excludes zero in 1 of 3 seeds at 1\% and in none at 10\%. A plausible reading is that the pretrained first layers were fitted to RGB, so the other bands add little, while the scratch CNN learns its own first layer; we did not test this. For all three models the gain is larger at 1\% than at 100\% of the labels, but it does not grow steadily as labels become scarcer: the scratch CNN gains more at 10\% than at 1\%, and ViT-Small gains slightly more at 100\% (0.15 points) than at 10\% (0.12). Per class, for ViT-Small at 1\% labels (Table~\ref{tab:per_class_f1_1pct}), the clearest gains are River (84.8 to 91.4\%) and SeaLake (95.3 to 97.7\%), and Residential and Highway lose 1.7 and 1.2 points; most class differences are within one to two standard deviations over seeds.

\subsection{Cross-check with an independent pipeline}\label{sec:crosscheck}
A teammate, Ayush Kumar, built an independent pipeline for the same task (the folder \texttt{ayush\_pipeline/} of the repository; a different split, recipe and seeds; accuracy only). Table~\ref{tab:crosscheck} compares the RGB accuracy of the same three models. The pretrained models agree within 0.6 points in all six cells and the scratch CNN within 1.0 point (ours is 0.8 to 1.0 points lower), so the two pipelines give the same label-efficiency picture. His pipeline also trained ResNet-50 on the three RGB bands cut from the 13-band GeoTIFF files, which separates the effect of the input source from the effect of the extra bands (Table~\ref{tab:band_control}). At 1\% of the labels, RGB from the GeoTIFF files scores 1.0 point above the JPEG RGB, and adding all 13 bands changes the score by $-0.2$ points on top of that. Our ResNet-50 gain of 0.8 points (accuracy) over JPEG RGB at 1\% is therefore explained by the input source and not by the extra bands. At 10\% the two pipelines differ more (our gain is 1.0 point, his 0.0 over JPEG RGB), and at 100\% every effect is 0.4 points or less. For ResNet-50, then, our RGB-versus-13-band gain should not be attributed to the extra bands. The scratch CNN's gain of 3.5 to 4.2 points is much larger than this source effect, but we did not measure the source effect for that model.

His band ablation (Table~\ref{tab:band_ablation}) feeds subsets of the 13-band files to the same network. Relative to RGB from the GeoTIFF files, the near-infrared band B08 adds 2.0 points at 1\% of the labels, and the red-edge set (B05--B07, B8A) and the short-wave-infrared set (B11, B12) add 2.4 and 2.8 points; these two differ from B08 alone by less than the seed-to-seed spread. From 5\% of the labels on, the near-infrared, red-edge and SWIR sets stay within 0.4 points of RGB, while all 13 bands are slightly worse at 1 to 10\% of the labels ($-0.2$ to $-0.7$ points). The useful extra information is therefore near-infrared, and it matters mostly when labels are very scarce.

\input{\tabdir/crosscheck.tex}
\input{\tabdir/band_control.tex}
\input{\tabdir/band_ablation.tex}

\subsection{Pseudo-labelling}
\begin{figure}[t]
\centering
\includegraphics[width=\columnwidth]{\figdir/pseudo_label_gain.png}
\caption{Change in test macro-F1 of each pseudo-label student against its own teacher (same model, input, fraction and seed); bars are means, dots single seeds. Students help most at 1\% of the labels.}
\label{fig:pl}
\end{figure}
\input{\tabdir/semi_supervised.tex}
Pseudo-labelling helps when labels are scarcest (Fig.~\ref{fig:pl}, Tables~\ref{tab:semi_supervised} and~\ref{tab:pl_details}). ResNet-50 students beat their teachers by 2.4 points on RGB and 2.3 on 13 bands at 1\% of the labels (individual students: +2.2 to +2.4), and by 0.2 points on both at 10\% (+0.1 to +0.3), in both seeds each. A pseudo-labelled ResNet-50 at 1\% (RGB 90.9\%, 13 bands 92.0\%) reaches the supervised ViT-Small at 1\% (90.7\% and 91.2\%). The scratch CNN student, one seed, gains 1.8 points. The threshold shows why 10\% teachers add little: at 10\% labels, ResNet-50 teachers reach 95\% validation precision already at the lowest allowed $\tau=0.50$ and keep 99.7--99.8\% of the pool, so the students see mostly images the teacher already classifies correctly. At 1\% labels, $\tau$ was 0.87--0.94 with validation precision 0.950--0.952, keeping 85--91\% of the pool. The scratch CNN teacher at 1\% never reached 95\% precision ($\tau=0.99$, precision 0.893) and kept 64\% of the pool, so its pseudo-labels are noisier. We did not keep the accuracy of the pseudo-labels against the hidden pool labels, so we report validation precision only. A check on a deliberately weak teacher (outside these runs) showed that a global $\tau$ can leave whole classes without pseudo-labels; the class-balanced variant that addresses this was not run.

\input{\tabdir/pl_details.tex}

\subsection{Error analysis}
\begin{figure}[t]
\centering
\includegraphics[width=\columnwidth]{\figdir/per_class_f1_1pct.png}
\caption{Per-class test F1 of ViT-Small at 1\% of the labels (mean over three seeds), RGB and 13 bands.}
\label{fig:perclass}
\end{figure}
\begin{figure*}[t]
\centering
\includegraphics[width=0.85\textwidth]{\figdir/confusion_1pct.png}
\caption{Row-normalised confusion matrices of ViT-Small at 1\% labels, RGB and 13 bands, mean of three seeds. The confusions are mostly among vegetation classes (Pasture, PermanentCrop, HerbaceousVegetation, AnnualCrop).}
\label{fig:confusion}
\end{figure*}
At 1\% of the labels, the largest confusions of ViT-Small are among vegetation classes: Pasture predicted as PermanentCrop (6.0--6.7\% of Pasture images), Pasture as HerbaceousVegetation (4.7--5.1\%), HerbaceousVegetation as PermanentCrop (4.4--4.9\%), and AnnualCrop as PermanentCrop (4.8\%, 13 bands). On RGB the largest single confusion is River predicted as Highway (6.7\%), which falls to 4\% on 13 bands, in line with River's gain in F1. The vegetation classes stay at 86--93\% F1 with either input, so the extra bands do not remove the main errors in this setting.
\input{\tabdir/per_class_f1_1pct.tex}

With 10\% of the labels (Fig.~\ref{fig:confusion10}, Table~\ref{tab:per_class_f1_10pct}) the picture changes. Every class of ViT-Small reaches at least 94.5\% F1, and the River gain from 13 bands has disappeared (97.9\% with both inputs), so it was a small-sample effect. The largest remaining confusions are AnnualCrop predicted as PermanentCrop (2.1--2.8\%) and PermanentCrop predicted as HerbaceousVegetation (1.4\% on 13 bands, 2.8\% on RGB). PermanentCrop is the class with the lowest F1 on RGB (94.5\%) and gains 1.2 points on 13 bands, SeaLake gains 1.0 point (98.6 to 99.6\%), and Pasture loses 1.1 points. These per-class differences are small compared with their standard deviations over three seeds, so we do not interpret them one by one.

\begin{figure*}[t]
\centering
\includegraphics[width=0.85\textwidth]{\figdir/confusion_10pct.png}
\caption{Row-normalised confusion matrices of ViT-Small at 10\% labels, RGB and 13 bands, mean of three seeds. Most off-diagonal cells are below 2\%.}
\label{fig:confusion10}
\end{figure*}
\input{\tabdir/per_class_f1_10pct.tex}

\section{Discussion}\label{sec:discussion}
\subsection{Answers to the three questions}
\emph{How much does accuracy fall with fewer labels?} Little for pretrained models and a lot for a model trained from scratch: from 100\% to 1\% of the labels, the pretrained models lose 8--10 macro-F1 points and the scratch CNN 21--23; at 10\% the pretrained models are within 1.3--2.2 points of their 100\% score and the scratch CNN on RGB is 5.7 points below.

\emph{Do the extra bands help?} Yes, but how much depends on the model. The gain is 3.5--4.2 points for the scratch CNN, about 0.9 for ResNet-50, and 0.5 or less for ViT-Small, where at 1\% it is within the spread between seeds. The sign is the same in every seed for the scratch CNN and ResNet-50. For pretrained models, a large part of the benefit of more labels comes from ImageNet pretraining, which the extra bands cannot replace.

\emph{Does pseudo-labelling help?} At 1\% of the labels, yes: about +2.3 points for ResNet-50, enough to match the supervised ViT-Small. At 10\% the gain is 0.2 points, too small to matter. The evidence is two seeds for ResNet-50 and one for the scratch CNN, and the students use more images and steps than their teachers, so some of the gain may come from extra compute.

\subsection{Practical recommendations}
The results suggest four rules for a similar problem. First, when only a few hundred labels are available, start from a pretrained model: the gap to a model trained from scratch is 11--16 points at 1\% of the labels, and the pretrained transformer was the best choice at 1\% and 10\%. Second, when the labelled set is large, the choice of model matters little (all models are within 1.6 points at 100\%), so cost can decide: the scratch CNN trains in 22 minutes against 42 for ResNet-50. Third, use the extra bands if the model is trained from scratch, where they gave 3.5--4.2 points, but do not expect more than one point for a pretrained RGB network without a better first-layer initialisation (and for ResNet-50 even that point is explained by the input source, Section~\ref{sec:crosscheck}). Fourth, pseudo-labelling is worth its roughly 1.5-fold training cost at 1\% of the labels, where it lifted ResNet-50 to the level of the supervised ViT-Small, and not at 10\%, where the teacher is already accurate on almost the whole pool.

\subsection{Threats to validity}
\emph{Internal.} The learning rate was chosen on one cell and reused, so a model whose best rate differs by fraction or input may be under-tuned; the ViT-Small rate was not tuned at all. Few seeds (three, two or one) leave wide uncertainty, so small differences, such as those between the pretrained models at 100\%, should not be read as rankings. \emph{Construct.} Macro-F1 on EuroSAT is close to saturation at 100\% of the labels, so differences there carry little information. \emph{External.} We use one dataset with one split, $64\times64$ patches upsampled to $224\times224$, and models pretrained on natural images. Results may differ on larger patches, on other regions, or with models pretrained on satellite data. \emph{Reporting.} Most numbers come from printed logs, which means no bootstrap intervals for most comparisons; the direction of effects is supported more strongly than their exact size.

\section{Limitations}\label{sec:limits}
\begin{itemize}
\item A random split can overestimate accuracy, because EuroSAT patches come from a limited set of Sentinel-2 tiles and neighbouring patches look alike. We follow the protocol of the original paper; a spatially disjoint split would be stricter but would no longer match the benchmark.
\item The RGB images are 8-bit compressed renderings, while the 13-band files hold raw 16-bit values, so part of any RGB-versus-13-band difference may come from quantisation and compression, not only from the extra bands; a teammate's control (Section~\ref{sec:crosscheck}) puts this effect at about 1 point for ResNet-50 at 1\% of the labels.
\item Pretrained weights come from RGB images. Extra bands start from copies of the RGB filters, and a model pretrained on multispectral data would avoid this limit.
\item The band B10 (cirrus) carries almost no surface information (mean about 12 digital numbers) and is kept so that all 13 bands are used; we did not run an ablation without it.
\item Students train on more images and for more steps than their teachers, so part of any pseudo-labelling gain may come from the extra compute. The teacher's weights were also selected on the validation set used to choose $\tau$, so its validation precision can be slightly optimistic; we report the validation precision at $\tau$ and could not report the accuracy of the pseudo-labels against the hidden pool labels.
\item Three seeds give wide uncertainty; we report per-seed differences and how many intervals exclude zero, not p-values.
\item The compute budget was a few free-tier T4 sessions, so the 100\% cells have one seed and no error bars, ViT-Small has no learning-rate sweep and no pseudo-labelling, and the optional experiments were skipped. A ViT-Small that is better tuned could change its ranking against ResNet-50.
\end{itemize}

\section{Future Work}\label{sec:future}
Five extensions follow from the limitations. (i)~Start from a model pretrained on 13-band Sentinel-2 data~\cite{wang2023ssl4eo,cong2022satmae}, which would show how much of the small gain for pretrained models comes from the RGB-based first layer. (ii)~Use a spatially disjoint split to test how much of the accuracy comes from neighbouring, similar patches. (iii)~Run the missing seeds (100\% cells, the remaining students) and keep the prediction files of every run, which would give bootstrap intervals and per-class results for all models. (iv)~Add a band ablation (13 bands without B10, and RGB plus near-infrared) to see which extra bands carry the gain for the scratch CNN. (v)~Replace the global threshold by a per-class or class-balanced selection, and compare with consistency-based methods such as FixMatch~\cite{sohn2020fixmatch}, on the same compute budget as the supervised baseline~\cite{oliver2018realistic}.

\section{Conclusion}
On EuroSAT, a pretrained model keeps 88.5--91.2\% macro-F1 with 1\% of the labels, where a scratch CNN reaches 74.6--78.1\%, so pretraining matters more than the amount of labelled data. The 13-band input adds 3.5--4.2 points for the scratch CNN but under one point for the pretrained models, and a control from an independent pipeline shows that for ResNet-50 this small gain comes from the input source (GeoTIFF versus JPEG), not from the extra bands. Pseudo-labelling adds about 2.3 points for ResNet-50 at 1\% of the labels and little at 10\%. Because of the time limit, the 100\% results have one seed, pseudo-labelling has two seeds, and most comparisons have no bootstrap intervals, so we state the direction of effects and not exact sizes.

\bibliographystyle{IEEEtran}
\bibliography{references}

\appendices
\section{All Runs}\label{app:runs}
Table~\ref{tab:all_runs} lists every finished run. Scores are those printed at the end of each run; the validation F1 belongs to the weights selected by early stopping. The minutes are the training time on a Tesla T4, including the final prediction passes.

\input{\tabdir/all_runs.tex}

\section{Design Decisions}\label{app:design}
Table~\ref{tab:design} lists the main choices, why each was made, and the weakness we know of.

\begin{table*}[t]
\centering
\caption{Main design decisions, with the reason and the known weakness.}
\label{tab:design}
\footnotesize
\begin{tabular}{r p{4.6cm} p{5.2cm} p{5.6cm}}
\toprule
\# & Decision & Reason & Known weakness \\
\midrule
1 & One fixed random, class-stratified 70/15/15 split (seed 42) & Matches the original EuroSAT protocol; every run shares the same validation and test sets & Neighbouring patches look alike, so scores can be optimistic for new regions \\
2 & Nested labelled subsets (1\% in 10\% in 100\%) with 3 seeds & Adding labels never removes images already in use, so curves are easier to compare & Only 3 seeds, so uncertainty is large at 1\% \\
3 & Normalisation statistics from the training split only & Nothing from validation or test leaks into preprocessing & The statistics include images that count as unlabelled at 1\% and 10\% (no labels are used) \\
4 & Best weights by validation macro-F1 with early stopping; test set predicted once & The test set never influences any choice & Validation scores wobble at 1\% labels, so the selected epoch is noisy \\
5 & Equal-budget rule: $\max(\text{epochs}\times\text{one pass},\text{min steps})$ & Small label sets get enough steps to converge, with the same number of validation checks for every fraction & The 1\% and 10\% runs get more passes over the data than the 100\% runs \\
6 & 13-band first layer built from pretrained RGB filters (mean filter for other bands, scaled by 3/13) & Allows ImageNet weights with 13 bands and keeps the first-layer output size & ImageNet knows nothing about infrared; the RGB filters are weakened by the scaling \\
7 & Macro-F1 as the main metric (accuracy also reported) & Classes are mildly imbalanced and each class should count equally & Macro-F1 on 1\% runs is noisy \\
8 & Pseudo-label threshold $\tau$ chosen on validation (smallest $\tau$ with at least 95\% precision) & No test data is used, and as many pseudo-labels as possible are kept at a known quality & The teacher's weights were also selected on validation, so its precision can be slightly optimistic \\
9 & Students start from fresh pretrained or random weights, not from the teacher & The comparison with the teacher stays like for like & Students train on more images and steps; they run 10 epochs to limit extra compute \\
10 & Paired comparisons matched on seed; 95\% bootstrap interval over test images & Removes variation due to different labelled subsets & With 3 seeds, we report how many seeds agree and not p-values \\
11 & Time cuts: 100\% cells use one seed; no ViT-Small sweep; pseudo-labelling only at 1\% and 10\% & Needed to finish by the deadline & No standard deviation at 100\%; ViT-Small may be under-tuned \\
\bottomrule
\end{tabular}
\end{table*}

\section{Reproducibility}\label{app:repro}
The code, the logs of every run and all figures and tables are in the public repository \url{https://github.com/ISHANT57/p21}; the results in this report are the tag \texttt{results-v1}. The code is released under the MIT licence, and automated tests run on every change. Table~\ref{tab:repro} lists the steps, each a script in the \texttt{scripts} folder of the repository.

\begin{table}[t]
\centering
\caption{Steps to reproduce the results.}
\label{tab:repro}
\small
\begin{tabular}{rll}
\toprule
Step & Script & Purpose \\
\midrule
1 & \texttt{download\_data.py} & get both versions, verify hashes \\
2 & \texttt{make\_splits.py} & split and nested subsets \\
3 & \texttt{build\_cache.py} & pack images into arrays \\
4 & \texttt{run\_grid.py} & train the grid of runs \\
5 & \texttt{evaluate.py} & metrics from predictions \\
6 & \texttt{parse\_colab\_logs.py} & results from run logs \\
7 & \texttt{eval\_from\_logs.py} & tables from run logs \\
8 & \texttt{make\_figures.py} & report figures, tables \\
9 & \texttt{make\_report\_tables.py} & extra report tables \\
\bottomrule
\end{tabular}
\end{table}

Every run stores its configuration, the library versions, the GPU, the git commit and its predictions next to its metrics, so a number in this report can be traced to a run. The fixed seeds control the labelled subsets, the initialisation, the shuffling and the augmentation. The split files carry hashes, and training stops if a file has changed, so no result can silently come from another split.

\end{document}
